\documentclass[11pt]{article}
\usepackage[margin=1in]{geometry}
\usepackage[T1]{fontenc}
\usepackage{lmodern,amsmath,amssymb,amsthm,mathtools,hyperref}
\setlength{\emergencystretch}{3em}
\newtheorem{theorem}{Theorem}[section]
\newtheorem{lemma}[theorem]{Lemma}
\newtheorem{corollary}[theorem]{Corollary}
\newtheorem{proposition}[theorem]{Proposition}
\newcommand{\E}{\mathbb E}
\newcommand{\tr}{\operatorname{tr}}
\title{Constructive phase rounding, completion geometry,\\and obstructions to greedy repair}
\author{Joshua Mark\thanks{ORCID: \href{https://orcid.org/0009-0003-5905-2382}{0009-0003-5905-2382}. AI tools assisted with derivation, independent review, computation and exposition.}}
\date{20 November 2026 --- research draft; priority not established}
\begin{document}\maketitle

\begin{abstract}
We examine constructive steps and obstructions in four Hadamard
formulations. An explicit dephasing algorithm rounds sufficiently
small fourth-power dispersion into an exact quaternary Hadamard
matrix, including a two-objective certificate for approximately
orthogonal inputs. Real zero-supported completion of a two-zero
weighing matrix is equivalent to sign completion, with at most six
sign solutions beyond order four. The spectral branch derives exact
entry-move formulas and finite traps, then lifts these obstructions
to arbitrarily large known Hadamard orders. A rank-two residual
projection links two-column complex completion to qubit geometry;
classical real completion then implies automatic quaternary completion
when only two columns are missing.
All-order existence and bibliographic priority remain unresolved.
\end{abstract}
The conjecture asks for a real sign matrix $H$ with $HH^T=nI_n$
at every positive order divisible by four. Local certificates below
are not assumed to be attainable at every such order.
\section{A constructive phase-rounding certificate}
Let $F$ be a complex Hadamard matrix of order $n$, with entries of
modulus one and $FF^*=nI$. Set $M=F^{\circ4}$ and
\[
 D_4(F)=n^4-\|MM^*\|_F^2.
\]
The nonnegative deficit $D_4$ is invariant under row/column phases
and permutations. The previous pass proved an obstruction when
$\rho_4=\|M\|_{\rm op}/n$ is bounded away from one. Here we prove a
positive, dimension-dependent rounding statement.

\begin{theorem}[Finite dephasing and exact quaternary rounding]
\label{thm:phase-rounding}
If $D_4(F)<n$, or $D_4(F)<2n$ when $n$ is even,
row and column phases convert $F$ into
a matrix whose entrywise nearest-fourth-root rounding is a
quaternary Hadamard matrix of order $n$.
Only $n^2$ explicit dephasings need be tested.
In particular it suffices that
\[
 1-\rho_4(F)^4<
 \begin{cases}2/n^3,& n\text{ even},\\1/n^3,&n\text{ odd}.\end{cases}
\]
This sufficient deficit is of order $n^{-3}$ in the last inequality;
no such estimate is established at every even order.
\end{theorem}
\begin{proof}
For each pivot $(k,\ell)$ define the dephased matrix
\[
 F^{k\ell}_{ij}
 =F_{ij}\overline{F_{kj}}\,
       \overline{F_{i\ell}}F_{k\ell}.
\]
These are row and column phase operations, so each matrix is still
Hadamard. Its entrywise fourth power is the corresponding four-cycle
cross ratio of $M$. Direct expansion and summation give the exact identity
\[
 \sum_{k,\ell}\|(F^{k\ell})^{\circ4}-J_n\|_F^2=2D_4(F).
\]
Indeed the sum of the real parts of those cross ratios is
$\sum_{i,k}|\sum_j M_{ij}\overline{M_{kj}}|^2=\|MM^*\|_F^2$.
Some pivot therefore has fourth-power squared error at most
$2D_4/n^2$.

For a unit complex number $z$, let $q\in\{1,-1,i,-i\}$ be nearest.
Write $z/q=e^{it}$ with $|t|\le\pi/4$. Then
\[
 |z-q|^2=4\sin^2(t/2)\le
 \tfrac14\,4\sin^2(2t)=\tfrac14|z^4-1|^2.
\]
For example $\sin(2t)=4\sin(t/2)\cos(t/2)\cos t$, whose absolute
ratio to $\sin(t/2)$ is at least two on this interval.
Thus the rounded matrix $Q$ at that pivot has
$\eta=\|F^{k\ell}-Q\|_F^2\le D_4/(2n^2)$.
Write $e_i$ for the difference between row $i$ of $Q$ and the
corresponding row of $F^{k\ell}$. For distinct rows, orthogonality
and the fact that both matrices have row norm $\sqrt n$ give
\[
 |(QQ^*)_{ij}|
 \le\sqrt n(\|e_i\|+\|e_j\|)
 \le\sqrt{2n\eta}.
\]
Every off-diagonal Gram entry is a Gaussian integer. A nonzero one
has modulus at least one, so $\eta<1/(2n)$ suffices to force zero.
For even $n$, its real and imaginary integer parts have even sum:
it is a sum of an even number of fourth roots. A nonzero such entry
has modulus at least $\sqrt2$, so $\eta<1/n$ suffices.
The diagonal Gram entries are already $n$. This proves exact
quaternary orthogonality under the stated deficit bounds. Finally
$\|MM^*\|_F^2\ge\|M\|_{\rm op}^4=n^4\rho_4^4$ proves the last claim.
\end{proof}

The proof is an algorithm: dephase at all pivots, choose one minimizing
the fourth-power error, round, and check the Gaussian-integer Gram.
Alternatively try every pivot and stop at an exact Gram certificate.
No target Hadamard matrix is needed as an input. Computing all pivot
errors naively uses $O(n^4)$ operations. Computing row inner products
of $M$ first permits the error scores to be obtained in $O(n^3)$
operations: the score at $(k,\ell)$ is
\[
 2n^2-2\operatorname{Re}\sum_i
   \overline{M_{i\ell}}M_{k\ell}
   \sum_j M_{ij}\overline{M_{kj}}.
\]
Rounding and checking one selected pivot use $O(n^3)$ additional
operations. Arithmetic complexity here counts exact scalar operations;
it is not a bit-complexity bound for arbitrary real inputs.

\begin{corollary}[Separation in the absence of a quaternary matrix]
If $\mathrm{BH}(n,4)$ does not exist, every complex Hadamard matrix
of that order satisfies
\[
 D_4(F)\ge
 \begin{cases}2n,&n\text{ even},\\n,&n\text{ odd},\end{cases}
 \qquad
 1-\rho_4(F)^4\ge
 \begin{cases}2/n^3,&n\text{ even},\\1/n^3,&n\text{ odd}.\end{cases}
\]
\end{corollary}
This is a contrapositive separation theorem, not a construction of a
point below the threshold. At odd $n>1$ quaternary existence is
impossible: a vanishing sum of an odd number of fourth roots has
unequal parity in its real and imaginary counts. At even $n$,
an actual threshold-crossing construction remains the central task.
The classical fourth-cycle invariants and Butson realification are
background; see Tadej--\.Zyczkowski,
\url{https://arxiv.org/abs/quant-ph/0512154}.
No bibliographic novelty claim is made for the rounding estimate.

\subsection{A certificate that does not assume exact orthogonality}
The same algorithm also accepts approximate inputs.
Let $F$ now be any $n$-square unit-modulus matrix, define $D_4$ as
before, and put $\beta=\max_{i\ne j}|(FF^*)_{ij}|$.
Define $\gamma_n=\sqrt2$ for even $n$ and $\gamma_n=1$ otherwise.
\begin{proposition}[Two-objective certificate]
If
\[
 \beta+\sqrt{D_4(F)/n}<\gamma_n,
\]
one of the pivot dephasings rounds to an exact $\mathrm{BH}(n,4)$.
\end{proposition}
\begin{proof}
The averaging identity and nearest-root estimate only use unit modulus.
Dephasing preserves the magnitudes of the original Gram entries.
For the selected pivot and its rounded $Q$, the same two-row
estimate, now retaining the original correlation, gives
\[
 |(QQ^*)_{ij}|\le\beta+\sqrt{2n\eta}
 \le\beta+\sqrt{D_4/n}<\gamma_n.
\]
Gaussian-integrality and parity again force zero.
\end{proof}
Thus a continuous search could minimize correlation and phase dispersion
jointly, then obtain a discrete exact certificate on crossing this
explicit boundary. This is not a convergence theorem. In particular a
real sign matrix already has $D_4=0$, but its even integer correlations
must already vanish to satisfy this sufficient condition. The certificate
does not turn arbitrary sign-matrix local minima into Hadamard matrices.
\section{Exactness of the real relaxation for two-zero weighing completion}
\label{sec:real-weighing-exactness}

The eight-state completion test admits a stronger interpretation.
For a weighing matrix with two zeros per row at an admissible order,
allowing arbitrary real entries in the holes does not enlarge the
class of matrices that can be completed. A real completion either
already supplies a sign state or produces a conference matrix from
which a sign completion can be recovered. We also classify the
continuous solution families and sharpen the eight-fill bound.

Throughout $4\mid n$, $n\ge4$, and $W$ is $W(n,n-2)$. Use the
previous pass's row and column normalization: all zero positions
are diagonal two-square blocks, and
\[
 F=\begin{pmatrix}1&1\\1&-1\end{pmatrix},\qquad
 A_{ij}=\tfrac12W_{ij}F\in\mathcal D\quad(i\ne j),
\]
where $\mathcal D$ is the eight-element signed-permutation group.
Put
\[
 I=\begin{pmatrix}1&0\\0&1\end{pmatrix},\quad
 J=\begin{pmatrix}0&1\\-1&0\end{pmatrix},\quad
 D=\begin{pmatrix}1&0\\0&-1\end{pmatrix},\quad
 S=\begin{pmatrix}0&1\\1&0\end{pmatrix}.
\]
Let $\mathcal R=\operatorname{span}_{\mathbb R}\{I,J\}$ and
$\mathcal F=\operatorname{span}_{\mathbb R}\{D,S\}$.
The orthogonal group $O(2)$ consists exactly of the unit coefficient
circles in these two planes.

\subsection{A linear kernel controls every real completion}
A real zero-supported completion has diagonal fill blocks
$X_i=R_iF$, $R_i\in O(2)$, and satisfies
\[
 A_{ij}R_j^{\mathsf T}+R_iA_{ji}^{\mathsf T}=0.
\]
For a root $R=R_0$, define
\[
 R_j=-A_{j0}R^{\mathsf T}A_{0j}\quad(j>0).
\]
Let $K\subseteq M_2(\mathbb R)$ be the homogeneous linear space
where all remaining cross equations hold. A root gives a real
completion exactly when $R\in K\cap O(2)$; it gives a Hadamard
filling exactly when $R\in K\cap\mathcal D$.

\begin{lemma}[The only possible component sections]
Before imposing any additional consequences of the weighing
property, each of $K\cap\mathcal R$ and $K\cap\mathcal F$ is zero,
the whole plane, a coordinate-axis line, or one of the two diagonal
lines in the displayed coefficient bases.
\end{lemma}
\begin{proof}
After multiplying the cross equation for $i,j>0$ on the left and
right by signed permutations, it becomes
\[
 R+L_{ij}R^{\mathsf T}Q_{ij}=0,
\]
where
\[
 L_{ij}=A_{0j}A_{ij}^{\mathsf T}A_{i0},\qquad
 Q_{ij}=A_{0i}A_{ji}^{\mathsf T}A_{j0}
\]
are signed permutations. The operator
$R\mapsto-L_{ij}R^{\mathsf T}Q_{ij}$ is Frobenius-orthogonal.
It either interchanges the two planes, forcing their pure-component
solutions to zero, or preserves them and acts on each coefficient
pair by a signed permutation. The fixed space of a two-dimensional
signed permutation is zero, a full plane, an axis line, or a
diagonal line. Intersecting these fixed spaces introduces no other
line directions. This proves the assertion.
\end{proof}

\subsection{A diagonal real state produces a conference completion}
The next argument invokes the classical conference normalization
theorem, not an unproved rounding principle. Geramita--Verner,
\emph{Orthogonal Designs with Zero Diagonal}, Theorem 2,
\url{https://doi.org/10.4153/CJM-1976-028-3}, restates the
Delsarte--Goethals--Seidel result: a conference matrix at order
$0\bmod4$ can be made skew by sign changes. Equivalently its
unique zero per row and column can be filled by a signed permutation
$Y$ so that $U+Y$ is Hadamard. Both fills $Y$ and $-Y$ work.

\begin{lemma}[Conference conversion]
If $K$ contains a diagonal line in either coefficient plane, then
that entire plane is contained in $K$.
\end{lemma}
\begin{proof}
Normalize a diagonal root to an orthogonal matrix $R$. Its entries
are all $\pm1/\sqrt2$. Every propagated $R_i$ has the same entry
form, so $R_iF$ has exactly one nonzero $\pm\sqrt2$ in each row
and column. Thus the real fill is $X=\sqrt2T$, where $T$ is a
block diagonal signed permutation supported on half the zero
positions of $W$.

The completion equations give
$WT^{\mathsf T}+TW^{\mathsf T}=0$. Hence
$C=WT^{\mathsf T}$ is skew, has diagonal two-square blocks zero,
and satisfies $CC^{\mathsf T}=(n-2)I$. The integer matrix
\[
 U=C+I_n
\]
has entries $0,\pm1$ and $UU^{\mathsf T}=(n-1)I$: it is a
conference matrix, with its single remaining zero in each row at
the off-diagonal position of the corresponding coordinate pair.
Apply the classical conference theorem, first permuting its zero
matching onto the diagonal if needed. It supplies a signed
permutation $Y$ on these remaining zeros such that both $U\pm Y$
are Hadamard. The paired diagonal blocks of $U$ are $I_2$, so
orthogonality forces $Y_i^{\mathsf T}=-Y_i$, that is $Y_i=\pm J$.

Undoing $T$ gives the two sign completions
\[
 H_\pm=(U\pm Y)T=W+(I_n\pm Y)T.
\]
The matrices $T$ and $YT$ have complementary supports inside each
original zero block. Thus each hole receives one sign and every
nonzero entry of $W$ remains unchanged. At the root the two fill
states are
\[
 R_\pm=\frac{(I_2\pm Y_0)T_0F}{2}.
\]
They lie in the same orientation plane as the original diagonal
root $T_0F/\sqrt2$, and are linearly independent because $Y_0=\pm J$
has no real eigenvector. Both belong to $K$. Their span is therefore
the entire component plane, as asserted.
\end{proof}

\begin{theorem}[Exact real feasibility and solution topology]
A real matrix $H$ with $HH^{\mathsf T}=nI_n$, obtained by changing
only zeros of $W$, exists if and only if a zero-preserving Hadamard
filling exists. More precisely each component section
$K\cap\mathcal R$, $K\cap\mathcal F$ is only zero, a coordinate
axis, or the full plane. Each axis contributes two isolated
completions, both sign-valued. Each full plane contributes a full
circle of real completions containing exactly four sign fillings.
There are no isolated nonsign real completions.
\end{theorem}
\begin{proof}
The preceding lemmas exclude an isolated diagonal-line section.
Every nonzero permitted section intersects $O(2)$ and contains
signed-permutation roots. Conversely a real orthogonal root lies
in one of these sections. The counts follow by intersecting an
axis or a plane with its unit coefficient circle, and the root
propagation is injective. A nonsign real root can therefore occur
only on one of the full circles, not as an isolated solution.
\end{proof}

This is an exactness theorem for this specific real feasibility
relaxation. It does not claim that arbitrary approximate real
fillings can be stably rounded, or that real feasibility always
holds. The unfillable order-twelve example has both sections zero.

\subsection{At most six fillings beyond order four}
\begin{proposition}[A sharper cardinality bound]
For $n\ge8$ the number of zero-preserving Hadamard fillings belongs
to $\{0,2,4,6\}$ and is at most six. The bound six is attained.
At order four there are exactly eight fillings.
\end{proposition}
\begin{proof}
Each component contributes zero, two or four sign roots. If both
planes were full, $K$ would be all of $M_2(\mathbb R)$. At $n\ge8$
there is a block triangle and hence at least one equation
$R+LR^{\mathsf T}Q=0$. It cannot hold identically. Substituting
$R=I$ would give $Q=-L^{\mathsf T}$ and thus
$R=LR^{\mathsf T}L^{\mathsf T}$ for every matrix. This map reverses
products; if it were the identity, all two-square matrices would
commute, a contradiction. Thus $K$ is proper and eight is impossible.
The skew order-four conference tensor with $F$ has one full plane
and one axis, attaining six at order eight. For $n=4$ there are
only two blocks and no triangle constraints; every one of the
eight roots works. The displayed set gives necessary possibilities;
no assertion that every listed value occurs at every order is made.
\end{proof}

\subsection{Exact checks and a nonsign circle point}
The verifier checks the real kernel sections for the four basic
cases: order four has two full planes and eight fillings; the
skew-conference tensor at order eight has an axis and a full plane,
giving six; the symmetric-conference tensor at order twelve has
one axis and two fillings; the inconsistent order-twelve example
has no real or sign completion. It also exhausts all 256 normalized
$W(8,6)$ matrices with standard zero support and first-row/sign
normalizations. Every case has one axis and one full plane and
exactly six sign fillings.

For a concrete nonsign real completion, the order-eight full
reflection plane contains the rational root $(3D+4S)/5$.
Its resulting completion has exact Gram $8I$ and nonsign hole
entries. The converter is separately tested at diagonal states:
it constructs the two conference fills and checks both resulting
Hadamard matrices and the independence of their root states using
only integer and rational arithmetic. No isolated diagonal-line
example was found; the conversion theorem explains why none can
occur. No novelty or all-order Hadamard existence claim is made.
\section{Exact entry-flip calculus and persistent discrete traps}
\label{sec:spectral-flip-traps}
Put $F(H)=\|HH^T-nI\|_F^2$ for an order-$n$ sign matrix.
This discrete energy landscape differs from optimization over the real
orthogonal group. Squared-Gram spin formulations and nonglobal local
minima are established background; see the primary article
\emph{Quantum computing formulation of some classical Hadamard matrix
searching methods and its implementation on a quantum computer},
\url{https://doi.org/10.1038/s41598-021-03586-0}.
The earlier primary article \emph{Finding a Hadamard Matrix by Simulated
Quantum Annealing} (2018), \url{https://www.mdpi.com/1099-4300/20/2/141},
also studies an energy-search formulation.
The explicit certificates and tensor bounds below are derived directly;
novelty is not asserted.

\begin{lemma}[Exact one- and two-entry changes]\label{lem:spectral-flip-calculus}
Let $G=HH^T$, $K=H^TH$. Flipping entry $(i,a)$ changes energy by
\[
 d_{ia}=8\big[(2n-1)-h_{ia}(GH)_{ia}\big].
\]
For two distinct flipped cells the change is the sum of their individual
changes plus the interaction
\[
 \begin{cases}
 16(h_{ia}h_{ib}K_{ab}-1),&\text{same row }i,\ a\ne b,\\
 16(h_{ia}h_{ja}G_{ij}-1),&\text{same column }a,\ i\ne j,\\
 16h_{ia}h_{ib}h_{ja}h_{jb},&i\ne j,\ a\ne b
 \quad\text{for cells }(i,a),(j,b).
 \end{cases}
\]
For even $n$, every single-entry energy change is eight times an odd
integer, so there are no neutral single-entry moves.
\end{lemma}
\begin{proof}
Only row-$i$ correlations change in a single flip:
$G_{ij}\mapsto G_{ij}-2h_{ia}h_{ja}$ for $j\ne i$.
Summing the changes of their squared values, twice for symmetry, gives
$8[(n-1)-h_{ia}\sum_{j\ne i}h_{ja}G_{ij}]$, the displayed formula.
Alternatively expand
\[
 F(H)=n^2(n-1)+4\sum_{i<j,\ a<b}h_{ia}h_{ib}h_{ja}h_{jb}.
\]
The two-flip interaction is sixteen times the sum of rectangle products
containing both cells. This gives the three cases. For even $n$ all
entries of $G$ are even, so $(GH)_{ia}$ is even while $2n-1$ is odd.
\end{proof}

\begin{theorem}[An exact small trap and its optimal escape barrier]
\label{thm:spectral-eight-trap}
Let $A$ be the $8\times8$ sign matrix whose negative-entry bitmasks in
successive rows are
\[
 (84,44,30,230,133,73,111,51),
\]
where bit $j$ denotes column $j$, starting at zero. Then $F(A)=152$,
every single-entry flip increases energy by at least eight, and every
two-entry flip has nonnegative energy change. A three-entry flip lowers
energy to $128$. The minimum Hamming distance from $A$ to any
strictly lower-energy matrix is exactly three. The least energy barrier
along single-entry paths to a lower-energy matrix is exactly eight.
\end{theorem}
\begin{proof}[Exact finite certificate]
The matrix of single-flip changes divided by eight is
\[
 \begin{pmatrix}
 1&5&5&9&5&1&9&9\\
 5&9&1&1&5&9&9&5\\
 1&9&5&5&1&9&5&1\\
 5&5&1&5&5&1&5&1\\
 9&1&9&5&5&5&9&1\\
 1&1&5&1&5&5&5&5\\
 9&5&9&1&1&1&5&5\\
 5&1&1&9&1&5&5&9
 \end{pmatrix}.
\]
The checker evaluates every one of the $\binom{64}{2}=2016$ two-cell
changes using both direct Gram energy and the interaction formula;
their minimum is zero. It also evaluates all $\binom{64}{3}=41664$
three-cell changes, whose minimum is $-24$. Explicitly, flipping cells
with flattened indices $0,5,60$ gives energies
$152\to160\to152\to128$. The first step of any single-entry path
costs at least eight, so this path attains the minimum possible barrier.
The one- and two-flip checks prove the claimed distance-three minimum.
All calculations are exact integer operations, not a heuristic search
certificate. The original randomized hill search merely found the
matrix; none of the assertions depends on its search history.
\end{proof}

\begin{theorem}[Infinite strict two-flip traps]
\label{thm:spectral-tensor-discrete-traps}
For any sign Hadamard matrix $B$ of order $m\ge2$, the matrix
$A\otimes B$ has positive energy $152m^3$ and strictly increases energy
under every change of one or two entries. The exact minimum increases are
\[
 \min\Delta_1=16m-8,\qquad \min\Delta_2=32(m-1).
\]
In particular this gives infinitely many explicit nonglobal strict
one- and two-entry local minima at known orders $8\cdot2^r$, $r\ge1$.
\end{theorem}
\begin{proof}
Let $Q=AA^T$. The lift has row Gram $Q\otimes mI_m$, and
$HH^TH=(Q A)\otimes mB$. If a base cell has single-flip change $d$,
its lifted cells have change $md+8(m-1)$. Since the minimum base change
is eight, the exact lifted minimum is $16m-8$.
Energy scales as
$\|(Q-8I_8)\otimes mI_m\|_F^2=m^3F(A)$.

For two cells sharing neither row nor column, the interaction is at
least $-16$, so their total increase is at least $32m-32$.
For two cells in the same lifted row, if their Hadamard column indices
differ, column-Gram orthogonality again makes the interaction $-16$.
If those indices agree, their base columns differ and the exact lifted
pair change is $m\Delta_{2,\mathrm{base}}+32(m-1)$, nonnegative in its
first term. The same argument applies to a shared lifted column.
Thus every pair costs at least $32(m-1)$. Equality is attained by lifting
the neutral base pair $(0,0),(0,5)$ with the same Hadamard row and column
indices. Positive energy proves nonglobality, since tensoring Walsh order eight with $B$ gives an order-$8m$
Hadamard matrix.
\end{proof}

For completeness, the checker exhausts all $512$ normalized order-four
sign matrices. Only the six Hadamard matrices are single-entry local
minima. Row and column sign normalization preserves energy and the
entry-flip graph, so order eight is the first admissible order at which
this greedy-descent obstruction occurs. The example does not rule out
nonlocal methods: its base has a certified three-flip escape. It does
rule out a proof that arbitrary one- or two-entry improving moves must
always reach a Hadamard matrix. No all-order construction is supplied.
\section{A positive codimension-two bridge through qubit geometry}
Let $V\in\mathbb C^{n\times(n-2)}$ have unit-modulus entries and
$V^*V=nI_{n-2}$. Its residual Gram matrix
\[
 K=nI_n-VV^*
\]
is positive semidefinite of rank two, with diagonal entries two.
It has a factorization $K=RR^*$ with $R\in\mathbb C^{n\times2}$,
$R^*R=nI_2$, and every row of $R$ has squared norm two.
Define the real matrix
\[
 B_{ij}=\frac{|K_{ij}|^2}{2}-1.
\]

\begin{theorem}[An exact rank test for a two-column flat completion]
The matrix $B$ is a Gram matrix of unit vectors in $\mathbb R^3$.
The partial complex Hadamard matrix $V$ can be completed by two
unit-modulus columns if and only if $\operatorname{rank}B\le2$.
\end{theorem}
\begin{proof}
For each row $r_i$ of $R$, let $u_i=r_i^T/\sqrt2$ and form its unit
Bloch vector $b_i\in\mathbb R^3$. The elementary two-dimensional
identity $b_i\cdot b_j=2|u_i^*u_j|^2-1$ gives $B_{ij}$.
All other residual factors are obtained by a common two-dimensional
unitary change of basis. Such a factor has both entry moduli one
in every row exactly when every corresponding qubit has equal
coordinate probabilities, that is, every $b_i$ is perpendicular to
a fixed unit measurement axis. Such an axis exists exactly when
the Bloch vectors span dimension at most two. Explicitly, for a
unit normal $s$, diagonalize the traceless Hermitian matrix
$s\cdot\sigma$, whose eigenvalues are $\pm1$. Zero expectation in
each $u_i$ forces its two eigenbasis probabilities to be $1/2$.
This constructs the required residual factor.
\end{proof}

\begin{theorem}[Quaternary input needs no extra phase freedom]
Suppose $n$ is even and $V$ has entries in $\mu_4$.
Then $B_{ij}\in\{-1,0,1\}$, and $\operatorname{rank}B\le3$.
Whenever $\operatorname{rank}B\le2$, the two completing columns
can also be chosen in $\mu_4$. Thus complex and quaternary
completability are equivalent for this codimension-two input.
\end{theorem}
\begin{proof}
The off-diagonal $K_{ij}$ are Gaussian integers with real and
imaginary parts of even sum, since they are sums of the even number
$n-2$ of fourth roots. Positive semidefiniteness gives $|K_{ij}|\le2$.
Thus their squared magnitudes are $0,2,4$, proving the assertion on
$B$. Distinct Bloch vectors are consequently parallel, antiparallel,
or orthogonal, and lie on at most three mutually orthogonal axes.

If the rank is at most two, rotate their plane to the equator.
After a common phase choice the residual rows have the form
$a_i(1,\delta_i)$ with $|a_i|=1$ and
$\delta_i\in\{1,-1,i,-i\}$. For two nonorthogonal rows,
\[
 K_{ij}=a_i\overline{a_j}(1+\delta_i\overline{\delta_j}).
\]
If the two $\delta$ values agree, the Gaussian integer $K_{ij}$
has modulus two, so division by two is a fourth root. If the axes
differ, its modulus is $\sqrt2$, so division by $1+i$ or $1-i$
is again a fourth root. Therefore $a_i\overline{a_j}\in\mu_4$
on every nonorthogonality edge. For rank two this graph is connected:
states on one axis are nonorthogonal to every state on the other.
A global phase makes all $a_i$ fourth roots. In rank one the two
antipodal groups are disconnected but mutually orthogonal; choose
their common phases independently without changing $K$.
This also makes every $a_i$ a fourth root. The resulting residual
matrix has entries in $\mu_4$ and the same Gram $K$, and appending
it to $V$ gives the required quaternary matrix.
\end{proof}
The proof is an exact algorithm using only $K$: take a reference
Bloch axis from a row of $B$, choose an orthogonal second axis if
needed, and label their antipodes by $1,-1,i,-i$. Along nonzero
$K_{ij}$ edges solve the displayed phase ratio, starting one phase
at one in each connected component. All recovered phases are in
$\mu_4$. This avoids numerical extraction of a Bloch rotation.

\begin{theorem}[Automatic quaternary completion using a classical real input]
Every even-order quaternary partial Hadamard matrix missing two
columns has a quaternary completion.
This assertion uses the classical real codimension-four Hadamard
completion theorem as an input; it is not proved here from scratch.
\end{theorem}
\begin{proof}
The case $n=2$ is immediate. For $n\ge4$, realify $(1+i)V$ to a
real sign matrix of size $2n\times(2n-4)$ with column Gram $2nI$.
Classical real completion supplies four sign columns $Q$ with
\[
 QQ^T=2\mathcal R(K).
\]
On the other hand $\sqrt2\mathcal R(R)$ is a full-column-rank real
factor with that same Gram. The two factors differ by a right
orthogonal matrix, preserving the number of antipodal row directions.

If $B$ had rank three, the previous discrete inner-product argument
would give three orthogonal Bloch axes. Moreover $R^*R=nI_2$
implies $\sum_i b_i=0$, so both antipodes of each axis occur.
There would be at least six distinct complex lines represented by
the rows $r_i$. Realification contributes two different orthogonal
real row directions from each complex line; different complex lines
cannot share a nonzero real direction. Thus
$\sqrt2\mathcal R(R)$, and hence $Q$, would have at least twelve
antipodal real row classes. Four real sign coordinates have only
eight such classes. This is impossible. Therefore the Bloch rank is
at most two, and the preceding constructive theorem completes $V$
within $\mu_4$.
\end{proof}

The real codimension-four input is classical Hall--Verheiden
completion theory. Primary sources include E. Verheiden,
\emph{Integral and rational completions of combinatorial matrices},
J. Combinatorial Theory A 25 (1978), 267--276, and his dissertation
\url{https://thesis.caltech.edu/2899/}. The latter explicitly records
Hall's four-row theorem and its extension to seven. The criterion
and its two-dimensional quantum interpretation are presented as a
structural and algorithmic bridge; bibliographic novelty is not asserted.
The codimension-one statement is simpler: its residual has rank one
and diagonal one, and one of its columns directly supplies the
missing fourth-root column up to phase. Neither result supplies
the missing large partial frame at an unknown order.

The checker verifies the Bloch rank, phase-graph algorithm, and exact
Gaussian Gram identities at orders four, six and eight. A bounded
search of small quaternary partial frames found no rank-three input,
consistent with the classical-completion consequence. A three-axis
abstract Gaussian residual is also tested to distinguish the
residual geometry from a realizable quaternary partial frame.

\section{What the constructive results still need}
The phase-rounding criterion requires a starting matrix inside a
dimension-dependent basin. We have no all-order construction reaching
that basin. The low-codimension completion criteria require almost all
columns to have been constructed already. The weighing results preserve
specified entries and therefore solve a constrained completion task.
The spectral traps rule out a universal strict-descent algorithm based
only on one or two entry changes, but do not rule out larger moves,
temporary energy increases, or different objective functions.

The next constructive targets are an input-producing theorem for
phase rounding, a compatible all-order weighing construction, and
a global exchange operation with a proved escape mechanism.
No full-conjecture proof or previously unknown order is claimed.
No Lean edits, git commits, publication or submission were made.
\end{document}
